% Mathematical source: ../1-practice_activities/G11_T2_L1_C2_practice.tex
% User override: one worked, one guided, eight independent problems (1--10).
\makeatletter
\def\input@path{{../../../preamble/}}
\makeatother
\input{preamble.tex}
\usepackage{needspace}
\usepackage{varwidth}
\geometry{letterpaper,margin=0.25in}
\pagestyle{empty}
\linespread{1.04}
\setlength{\parskip}{4pt}
\raggedbottom
\newcommand{\stageheading}[1]{\par\addvspace{0.5\baselineskip}{\centering\large\bfseries #1\par}\nopagebreak[4]}

\newcommand{\problemtitle}[2]{\par\Needspace{4\baselineskip}\addvspace{.4\baselineskip}\noindent\textbf{#1. #2}\par\nopagebreak[4]}
\newcommand{\workspace}[1]{\par\vspace{#1}}
\newcommand{\task}[2]{\par\Needspace{3\baselineskip}\textbf{Task #1 -- #2}}
\newcommand{\choices}[4]{\par\begingroup\small\noindent A. #1\par B. #2\par C. #3\par D. #4\par\endgroup}
\newcommand{\InformationTask}{\textbf{Task 1 -- Identify and organize the information; calculate SE.} Use four columns: \textbf{Symbol}, \textbf{Concept}, \textbf{Type}, and \textbf{Meaning in context}, with one row per symbol and concept. Arrange three contiguous groups in this order: \textbf{Population parameters} ($\mu$, $\sigma$); \textbf{Sample statistics} ($\bar{x}$, $n$, $\mathrm{SE}$); \textbf{Procedure Settings} ($C$, $\alpha$, $\alpha/2$, $z^*$).}
\newcommand{\IntervalTask}{\textbf{Task 2 -- Calculate MoE and state the interval.} Use exactly five columns: \textbf{Symbol | Concept | Calculation | Result | Meaning in context}. Give the point estimate, margin of error, lower endpoint, upper endpoint, confidence interval, and full width separate rows. Put formulas and numerical substitutions in Calculation, final values or intervals with units in Result, and contextual explanations in Meaning in context. Group rows by confidence level and give the shared point estimate once. When rounding is required, distinguish unrounded and reported endpoints, intervals, and widths; calculate reported width from reported endpoints. Use $W=U-L$ and $W_{\rm rep}=U_{\rm rep}-L_{\rm rep}$.}
\newcommand{\CoverageTask}{\textbf{Task 3 -- Calculate observed coverage and answer the multiple-choice question.} Show the covering count, denominator, and coverage calculation separately for each procedure using the formula above. Select and justify one option.}

\newcommand{\InformationTable}[1]{%
    \par
    \begingroup
    \small
    \setlength{\tabcolsep}{4pt}
    \renewcommand{\arraystretch}{1.16}
    \noindent
    \begin{tabularx}{\linewidth}{@{}>{\raggedright\arraybackslash}p{0.07\linewidth}>{\raggedright\arraybackslash}p{0.18\linewidth}>{\raggedright\arraybackslash}p{0.19\linewidth}>{\raggedright\arraybackslash}X@{}}
        \toprule
        \textbf{Symbol} & \textbf{Concept} & \textbf{Type} & \textbf{Meaning in context} \\
        \midrule
        #1
        \bottomrule
    \end{tabularx}\par
    \vspace{2mm}
    \endgroup
}

% Each confidence-level block stays together; repeat headings in every block.
\newcommand{\EstimationTable}[1]{%
    \par
    \begingroup
    \setlength{\tabcolsep}{4pt}
    \renewcommand{\arraystretch}{1.18}
    \noindent
    \begin{tabularx}{\linewidth}{@{}>{\raggedright\arraybackslash}p{0.095\linewidth}>{\raggedright\arraybackslash}p{0.17\linewidth}>{\raggedright\arraybackslash}p{0.27\linewidth}>{\raggedright\arraybackslash}p{0.16\linewidth}>{\raggedright\arraybackslash}X@{}}
        \toprule
        \textbf{Symbol} & \textbf{Concept} & \textbf{Calculation} & \textbf{Result} & \textbf{Meaning in context} \\
        \midrule
        #1
        \bottomrule
    \end{tabularx}\par
    \vspace{2mm}
    \endgroup
}

\newcommand{\BlankInformationTable}{
\InformationTable{
$\mu$&Population mean&Population parameters&\\[23pt]
$\sigma$&Population standard deviation&&\\[18pt]\midrule
$\bar{x}$&Sample mean&Sample statistics&\\[18pt]
$n$&Sample size&&\\[18pt]
$\mathrm{SE}$&Standard error&&\\[27pt]\midrule
$C$&Confidence level&Procedure Settings&\\[18pt]
$\alpha$&Significance level&&\\[18pt]
$\alpha/2$&Tail area&&\\[18pt]
$z^*$&Critical value&&\\[18pt]}}
\begin{document}
{\centering{\Large\bfseries C2 -- Point and interval estimation\par}
Global Economics \textbar{} Grade 11 \textbar{} Term 2 \textbar{} Lesson 1\par}
\textbf{C2:} Compare point and interval estimation and interpret confidence through long-run coverage, distinguishing this from the probability that a particular calculated interval contains the parameter.

\stageheading{Equations and basic procedure}
Use independent random observations, known population SD $\sigma$, and a normal sampling distribution of $\bar{x}$ (exact for a normal population). The target $\mu$ is fixed and unknown; $\bar{x}$ is its point estimate. All problems below use these assumptions; simulation populations are normal with the stated known fixed means and the same $\sigma$ and $n$ as their corresponding original sample.

\textbf{Prerequisite:} Use the C1 reference table for population and sample concepts
and notation, including $\mu$, $\sigma$, $\bar{x}$, and $n$.
\begingroup
\small
\setlength{\tabcolsep}{4pt}
\renewcommand{\arraystretch}{1.12}
\noindent\begin{tabularx}{\linewidth}{@{}>{\raggedright\arraybackslash}p{.12\linewidth}>{\raggedright\arraybackslash}p{.23\linewidth}>{\raggedright\arraybackslash}X@{}}
\toprule
\textbf{Symbol} & \textbf{Concept} & \textbf{Formula} \\
\midrule
$C$ & Confidence level & \\
$z^*$ & Critical value & \\
$\mathrm{SE}$ & Standard error & $\mathrm{SE}=\frac{\sigma}{\sqrt n}$ \\
$\mathrm{MoE}$ & Margin of error & $\mathrm{MoE}=z^*\mathrm{SE}$ \\
$L$ & Lower endpoint & $L=\bar{x}-\mathrm{MoE}$ \\
$U$ & Upper endpoint & $U=\bar{x}+\mathrm{MoE}$ \\
$W$ & Full width & $W=U-L=2\mathrm{MoE}$ \\
$\alpha$ & Significance level & $\alpha=1-C$ \\
$\alpha/2$ & Each tail's area & $\frac{\alpha}{2}=\frac{1-C}{2}$ \\
$N_{\rm cover}(C)$ & Covering count & \\
$N(C)$ & Total interval count & \\
$\mathrm{Cov}(C)$ & Observed coverage & $\frac{N_{\rm cover}(C)}{N(C)}\times100\%$ \\
\bottomrule
\end{tabularx}\par
\endgroup
\begin{center}
\small
\textbf{Two-sided confidence levels and standard normal critical values}\par
\begin{tabularx}{.80\linewidth}{@{}>{\raggedright\arraybackslash}Xccccc@{}}\toprule
Confidence level $C$&80\%&90\%&95\%&98\%&99\%\\\midrule
Critical value $z^*$&1.282&1.645&1.960&2.326&2.576\\\bottomrule
\end{tabularx}\end{center}
A simulated interval covers the fixed simulation mean when $L\leq\mu_{\rm sim}\leq U$.
No sample SD $s$ is provided or needed.

\textbf{Precision and coverage.} Use the displayed critical values in this reference table as the calculation values. Keep all digits unless reporting-rounding is requested. A reported width is reported upper endpoint minus reported lower endpoint. Nominal confidence and simulation coverage refer to the unrounded procedures.


\stageheading{General assessment tasks}

\InformationTask

\IntervalTask

\CoverageTask

\clearpage
\stageheading{Worked example}
\vspace{-8mm}

\problemtitle{1}{Investigate the cycling interval}
A random sample of 49 cycling trips had an average duration of 40 minutes. The average duration of all trips in the population represented by this sample is unknown. The known population standard deviation of individual trip durations is 14 minutes. The confidence level for the report is 95\%.

In a separate illustrative simulation, five fresh independent samples of 49 trips were drawn from a normal population with known fixed mean $\mu_{\rm sim}=41$ minutes and population SD 14 minutes. Each sample uses the same 95\% procedure.
\begin{center}\begin{tabular}{@{}rrrr@{}}\toprule
Sample&Sample mean (min)&Lower endpoint (min)&Upper endpoint (min)\\\midrule
1&40&36.08&43.92\\
2&42&38.08&45.92\\
3&45&41.08&48.92\\
4&38&34.08&41.92\\
5&41&37.08&44.92\\\bottomrule
\end{tabular}\end{center}
\task{1}{Identify and organize the information; calculate SE.}
\task{2}{Calculate MoE and state the interval.}
\task{3}{Calculate observed coverage and answer the multiple-choice question.}
\choices{Four of five cover, so 80\% replaces 95\% as the procedure's long-run coverage.}
{Four of five cover. Observed coverage is 80\%; the 95\% procedure covers the fixed mean about 95\% of the time in the long run, and each calculated interval either covers or misses.}
{Four of five cover; the original calculated interval has a 95\% probability of containing its fixed unknown mean.}
{All five cover; 95\% confidence means that 95\% of individual trip durations lie in the original interval.}

\vspace{-3mm}
\subsubsection*{Task 1 solution}
\InformationTable{
$\mu$&Population mean&Population parameters&Fixed unknown mean duration of all cycling trips represented by the original sample, in minutes. The separate simulation mean is known: $\mu_{\rm sim}=41$ minutes.\\
$\sigma$&Population standard deviation&&$\sigma=14$ minutes, shared by the original and simulation populations; spread of individual trip durations.\\\midrule
$\bar{x}$&Sample mean&Sample statistics&$\bar{x}=40$ minutes is the original point estimate of $\mu$. Each fresh simulation sample has its own mean.\\
$n$&Sample size&&$n=49$ trips in the original sample and in each simulation sample.\\
$\mathrm{SE}$&Standard error&&$\mathrm{SE}=\sigma/\sqrt n=14/\sqrt{49}=2$ minutes. Spread of means across repeated independent samples of 49 trips, shared by these procedures.\\\midrule
$C$&Confidence level&Procedure Settings&95\% procedure: $C=0.95$, shared by original and simulated intervals.\\
$\alpha$&Significance level&&$\alpha=1-C=1-0.95=0.05$, total standard normal area outside the central region.\\
$\alpha/2$&Tail area&&$\alpha/2=0.05/2=0.025$ in each tail.\\
$z^*$&Critical value&&$z^*=1.960$, the standard normal multiplier for the 95\% procedure.\\}

\clearpage
\problemtitle{1}{Cycling interval -- solutions continued}
\subsubsection*{Task 2 solution}
\EstimationTable{
$\bar{x}$ & Point estimate & $\bar{x}=\dfrac{\sum x_i}{n}$; given sample mean & 40 minutes & Estimates the unknown population mean cycling-trip duration.\\[6pt]\midrule
\multicolumn{5}{@{}l}{\textbf{95\% confidence}}\\[3pt]
$\mathrm{MoE}$ & Margin of error & $z^*\mathrm{SE}=1.960(2)$ & 3.92 minutes & Distance from the estimated mean to either endpoint.\\[6pt]
$L$ & Lower endpoint & $\bar{x}-\mathrm{MoE}$\par $=40-3.92$ & 36.08 minutes & Lower bound for the plausible population mean duration.\\[6pt]
$U$ & Upper endpoint & $\bar{x}+\mathrm{MoE}$\par $=40+3.92$ & 43.92 minutes & Upper bound for the plausible population mean duration.\\[6pt]
$[L,U]$ & Confidence interval & $[\bar{x}-\mathrm{MoE},$\par $\bar{x}+\mathrm{MoE}]$\par $=[40-3.92,40+3.92]$ & $[36.08,43.92]$ minutes & We are 95\% confident that the population mean cycling-trip duration is in this range of plausible values under the model.\\[6pt]
$W$ & Full width & $U-L=43.92-36.08$\par $=2(3.92)$ & 7.84 minutes & Total span of the interval for the population mean duration.\\[6pt]}
Check the centre: $(36.08+43.92)/2=40$ minutes. The simulation intervals have the same 7.84-minute width, but their centres follow their fresh sample means. With the same sample and SE, higher confidence increases $z^*$, MoE, and width while preserving the centre, so the wider interval contains the narrower one. Increasing $n$ with other settings fixed reduces SE and width.

\subsubsection*{Task 3 solution}
The fixed simulation mean is 41 minutes. Samples 1, 2, 4, and 5 cover it because their endpoints satisfy $L\leq41\leq U$. Sample 3 misses: its lower endpoint, 41.08, exceeds 41.
\[
\text{95\% covering count}=4,\qquad
\text{95\% denominator}=5,\qquad
\text{observed coverage}=\frac45\times100\%=80\%.
\]
\textbf{Correct option B.} The observed 80\% is a finite-run result; over many fresh samples under the model, about 95\% of intervals made by this procedure cover the fixed population mean. Option A incorrectly replaces that long-run confidence with the observed rate. Option C assigns probability to an already calculated interval, whose fixed endpoints either contain the fixed unknown mean or do not. Option D miscounts sample 3 and confuses estimation of the population mean with the spread of individual trip durations.

\clearpage
\stageheading{Guided Example}
\problemtitle{2}{Two reports on weekly printing costs}
A random sample of 100 school offices in a district had an average weekly printing cost of 86 thousand COP. The average weekly cost for all offices in the district is unknown. The known population standard deviation of individual offices' weekly costs is 32 thousand COP. Two reports, at 90\% and 99\% confidence, concern this same sample.

In a separate, completed repeated-sampling study, 200 fresh independent samples of 100 offices were drawn from a normal population with a fixed known average weekly cost of 88 thousand COP and a population standard deviation of 32 thousand COP. Both unrounded confidence intervals were constructed from each sample. The 90\% intervals missed the known mean 23 times, and the 99\% intervals missed it 5 times.

\textbf{Reporting requirement.} Report endpoints to the nearest thousand COP.


\task{1}{Identify and organize the information; calculate SE.}
\InformationTable{
$\mu$&Population mean&Population parameters&Unknown district mean: \rule{2in}{.3pt}\par Known simulation mean: \rule{1.4in}{.3pt}\\[10pt]
$\sigma$&Population standard deviation&&Shared known SD: \rule{1.8in}{.3pt}\\[10pt]\midrule
$\bar{x}$&Sample mean&Sample statistics&Shared original point estimate: \rule{1.5in}{.3pt}\\[10pt]
$n$&Sample size&&Shared sample size: \rule{1.8in}{.3pt}\\[10pt]
$\mathrm{SE}$&Standard error&&$\mathrm{SE}=\sigma/\sqrt n=32/\sqrt{100}=$ \rule{.7in}{.3pt}\par Units and meaning: \rule{1.8in}{.3pt}\\[12pt]\midrule
$C$&Confidence level&Procedure Settings&90\%: $C=0.90$; 99\%: $C=$ \rule{.6in}{.3pt}\\[10pt]
$\alpha$&Significance level&&90\%: $\alpha=1-0.90=$ \rule{.5in}{.3pt}\par 99\%: $\alpha=1-$ \rule{.5in}{.3pt} $=$ \rule{.5in}{.3pt}\\[10pt]
$\alpha/2$&Tail area&&90\%: \rule{1in}{.3pt}; 99\%: \rule{1in}{.3pt}\\[10pt]
$z^*$&Critical value&&90\%: \rule{1in}{.3pt}; 99\%: \rule{1in}{.3pt}\\[10pt]}


\clearpage
\problemtitle{2}{Weekly printing costs -- continued}
\task{2}{Calculate MoE and state the interval.}
\EstimationTable{
$\bar{x}$ & Shared point estimate & Given sample mean & \rule{.7in}{.3pt}\par Units: \rule{.6in}{.3pt} & What district mean does this estimate?\\[18pt]\midrule
\multicolumn{5}{@{}l}{\textbf{90\% confidence}}\\[3pt]
$\mathrm{MoE}$ & Margin of error & $z^*\mathrm{SE}=1.645(\quad)$ & & Distance from centre to either bound.\\[20pt]
$L$ & Lower endpoint (unrounded) & $86-\mathrm{MoE}$ & & Lower bound for which mean?\\[20pt]
$U$ & Upper endpoint (unrounded) & $86+\mathrm{MoE}$ & & Upper bound for the same mean.\\[20pt]
$[L,U]$ & Confidence interval (unrounded) & $[\bar{x}-\mathrm{MoE},$\par $\bar{x}+\mathrm{MoE}]$ & & Interpret as plausible values for the unknown mean.\\[20pt]
$W$ & Full width (unrounded) & $U-L=2\mathrm{MoE}$ & & Interpret the total span.\\[20pt]
$L_{\rm rep}$ & Reported lower endpoint & Round $L$ to the nearest thousand COP. & & Explain the reported lower bound.\\[20pt]
$U_{\rm rep}$ & Reported upper endpoint & Round $U$ to the nearest thousand COP. & & Explain the reported upper bound.\\[20pt]
$[L_{\rm rep},U_{\rm rep}]$ & Reported confidence interval & $[L_{\rm rep},U_{\rm rep}]$ & & Interpret this reported range for the same mean.\\[20pt]
$W_{\rm rep}$ & Reported full width & $U_{\rm rep}-L_{\rm rep}$ & & Interpret the reported span.\\[20pt]
}
\clearpage
\problemtitle{2}{Weekly printing costs -- Task 2 continued}
\EstimationTable{
\multicolumn{5}{@{}l}{\textbf{99\% confidence}}\\[3pt]
$\mathrm{MoE}$ & Margin of error & $z^*\mathrm{SE}=2.576(\quad)$ & & Distance from centre to either bound.\\[20pt]
$L$ & Lower endpoint (unrounded) & $\bar{x}-\mathrm{MoE}$ & & Lower bound for which mean?\\[20pt]
$U$ & Upper endpoint (unrounded) & $\bar{x}+\mathrm{MoE}$ & & Upper bound for the same mean.\\[20pt]
$[L,U]$ & Confidence interval (unrounded) & $[\bar{x}-\mathrm{MoE},$\par $\bar{x}+\mathrm{MoE}]$ & & Interpret as plausible values for the unknown mean.\\[20pt]
$W$ & Full width (unrounded) & $U-L=2\mathrm{MoE}$ & & Interpret the total span.\\[20pt]
$L_{\rm rep}$ & Reported lower endpoint & Round $L$ to the nearest thousand COP. & & Explain the reported lower bound.\\[20pt]
$U_{\rm rep}$ & Reported upper endpoint & Round $U$ to the nearest thousand COP. & & Explain the reported upper bound.\\[20pt]
$[L_{\rm rep},U_{\rm rep}]$ & Reported confidence interval & $[L_{\rm rep},U_{\rm rep}]$ & & Interpret this reported range for the same mean.\\[20pt]
$W_{\rm rep}$ & Reported full width & $U_{\rm rep}-L_{\rm rep}$ & & Interpret the reported span.\\[20pt]
}
Explain why rounding can change the reported widths without changing the nominal confidence of the unrounded procedures.
\workspace{.65in}
\task{3}{Calculate observed coverage and answer the multiple-choice question.}
\par 90\%: covering count $=200-23=$ \rule{.6in}{.3pt}; denominator $=$ \rule{.5in}{.3pt}; coverage $=$ \rule{1.4in}{.3pt}.
\par 99\%: covering count $=200-5=$ \rule{.6in}{.3pt}; denominator $=$ \rule{.5in}{.3pt}; coverage $=$ \rule{1.4in}{.3pt}.
\choices{The 99\% interval contains 99\% of individual offices' printing costs.}
{Each finite simulation must cover at exactly its nominal confidence level.}
{The wider 99\% procedure covers the fixed mean about 99\% of the time in the long run; each calculated interval either covers or misses.}
{The reported 99\% interval has a 99\% probability of containing its fixed unknown mean.}
\workspace{.8in}


\clearpage
\stageheading{Independent Practice}
\problemtitle{3}{Repeated estimates of locker rental charges}
A random sample of 64 customers at a sports centre had an average weekly locker rental charge of 72 thousand COP. The average weekly charge for all customers at the centre is unknown. The known population standard deviation of individual customers' charges is 20 thousand COP. The confidence level for the report is 95\%.

In a separate, completed repeated-sampling study, five fresh independent samples of 64 customers were drawn from a normal population with a fixed known average charge of 70 thousand COP and the same population standard deviation as the original population. A 95\% confidence interval was constructed from each sample. Their endpoints appear below in thousand COP.

\begin{center}\begin{tabular}{@{}lll@{}}\toprule
Sample & Lower endpoint & Upper endpoint\\\midrule
1 & 64.1 & 73.9\\
2 & 67.1 & 76.9\\
3 & 70.1 & 79.9\\
4 & 63.1 & 72.9\\
5 & 65.1 & 74.9\\
\bottomrule\end{tabular}\end{center}


\task{1}{Identify and organize the information; calculate SE.}
\BlankInformationTable

\clearpage
\problemtitle{3}{Continued -- interval estimation and coverage}

\task{2}{Calculate MoE and state the interval.}
\EstimationTable{
$\bar{x}$ & Point estimate & & & \\[30pt]\midrule
\multicolumn{5}{@{}l}{\textbf{95\% confidence}}\\[3pt]
$\mathrm{MoE}$ & Margin of error & & & \\[30pt]
$L$ & Lower endpoint & & & \\[30pt]
$U$ & Upper endpoint & & & \\[30pt]
$[L,U]$ & Confidence interval & & & \\[30pt]
$W$ & Full width & & & \\[30pt]
}
\task{3}{Calculate observed coverage and answer the multiple-choice question.}

\choices{Four of five cover; 80\% is therefore the long-run coverage.}{Four of five cover; the procedure covers about 95\% in the long run, and this particular interval either covers or misses.}{Four of five cover; the original interval has a 95\% probability of containing its fixed mean.}{All five cover; the procedure covers about 95\% in the long run.}
\vfill


\clearpage
\problemtitle{4}{Repeated estimates of nursery delivery charges}
A random sample of 81 nursery orders in a specified month had an average delivery charge of 54 thousand COP. The average delivery charge for all nursery orders that month is unknown. The known population standard deviation of individual orders' charges is 18 thousand COP. The confidence level for the report is 95\%.

In a separate, completed repeated-sampling study, five fresh independent samples of 81 orders were drawn from a normal population with a fixed known average charge of 55 thousand COP and the same population standard deviation as the original population. A 95\% confidence interval was constructed from each sample. Their endpoints appear below in thousand COP.

\begin{center}\begin{tabular}{@{}lll@{}}\toprule
Sample & Lower endpoint & Upper endpoint\\\midrule
1 & 50.08 & 57.92\\
2 & 51.08 & 58.92\\
3 & 46.08 & 53.92\\
4 & 52.08 & 59.92\\
5 & 55.08 & 62.92\\
\bottomrule\end{tabular}\end{center}


\task{1}{Identify and organize the information; calculate SE.}
\BlankInformationTable

\clearpage
\problemtitle{4}{Continued -- interval estimation and coverage}

\task{2}{Calculate MoE and state the interval.}
\EstimationTable{
$\bar{x}$ & Point estimate & & & \\[30pt]\midrule
\multicolumn{5}{@{}l}{\textbf{95\% confidence}}\\[3pt]
$\mathrm{MoE}$ & Margin of error & & & \\[30pt]
$L$ & Lower endpoint & & & \\[30pt]
$U$ & Upper endpoint & & & \\[30pt]
$[L,U]$ & Confidence interval & & & \\[30pt]
$W$ & Full width & & & \\[30pt]
}
\task{3}{Calculate observed coverage and answer the multiple-choice question.}

\choices{Three of five cover; the calculated interval has a 95\% probability of containing the fixed mean.}{Five of five cover; confidence is about repeated intervals.}{Three of five cover; the procedure covers about 95\% in the long run, while a fixed calculated interval either covers or misses.}{Three of five cover; 60\% replaces 95\% as the long-run confidence.}
\vfill


\clearpage
\problemtitle{5}{Two reports on monthly pet-care spending}
A random sample of 144 registered households in a district had an average monthly pet-care expenditure of 145 thousand COP. The average monthly expenditure for all registered households in the district is unknown. The known population standard deviation of individual households' monthly spending is 42 thousand COP. Two reports, at 90\% and 99\% confidence, concern this same sample.

In a separate, completed repeated-sampling study, 150 fresh independent samples of 144 households were drawn from a normal population with a fixed known average expenditure of 148 thousand COP and the same population standard deviation as the original population. Both unrounded confidence intervals were constructed from each sample. The 90\% intervals missed the known mean 17 times, and the 99\% intervals missed it 3 times.

\textbf{Reporting requirement.} Round reported endpoints to the nearest thousand COP.


\task{1}{Identify and organize the information; calculate SE.}
\BlankInformationTable

\clearpage
\problemtitle{5}{Continued -- interval estimation and coverage}

\task{2}{Calculate MoE and state the interval.}
\EstimationTable{
$\bar{x}$ & Shared point estimate & & & \\[30pt]\midrule
\multicolumn{5}{@{}l}{\textbf{90\% confidence}}\\[3pt]
$\mathrm{MoE}$ & Margin of error & & & \\[30pt]
$L$ & Lower endpoint (unrounded) & & & \\[30pt]
$U$ & Upper endpoint (unrounded) & & & \\[30pt]
$[L,U]$ & Confidence interval (unrounded) & & & \\[30pt]
$W$ & Full width (unrounded) & & & \\[30pt]
$L_{\rm rep}$ & Reported lower endpoint & & & \\[30pt]
$U_{\rm rep}$ & Reported upper endpoint & & & \\[30pt]
$[L_{\rm rep},U_{\rm rep}]$ & Reported confidence interval & & & \\[30pt]
$W_{\rm rep}$ & Reported full width & & & \\[30pt]
}
\clearpage
\problemtitle{5}{Continued -- Task 2}
\EstimationTable{
\multicolumn{5}{@{}l}{\textbf{99\% confidence}}\\[3pt]
$\mathrm{MoE}$ & Margin of error & & & \\[30pt]
$L$ & Lower endpoint (unrounded) & & & \\[30pt]
$U$ & Upper endpoint (unrounded) & & & \\[30pt]
$[L,U]$ & Confidence interval (unrounded) & & & \\[30pt]
$W$ & Full width (unrounded) & & & \\[30pt]
$L_{\rm rep}$ & Reported lower endpoint & & & \\[30pt]
$U_{\rm rep}$ & Reported upper endpoint & & & \\[30pt]
$[L_{\rm rep},U_{\rm rep}]$ & Reported confidence interval & & & \\[30pt]
$W_{\rm rep}$ & Reported full width & & & \\[30pt]
}
\textbf{Compare the centre, widths, and nesting of the intervals.}
Explain how rounding affects the reported widths, while nominal confidence refers to the unrounded procedures.
\workspace{.65in}

\task{3}{Calculate observed coverage and answer the multiple-choice question.}

\choices{The wider 99\% interval contains 99\% of individual household spending amounts.}{Both finite coverage rates must equal their nominal confidence levels.}{The wider reported interval has a 99\% probability of containing the fixed mean.}{The 99\% unrounded procedure covers the fixed mean about 99\% of the time in the long run; rounding is a separate reporting step.}
\vfill


\clearpage
\problemtitle{6}{Two reports on weekly bakery supply costs}
A random sample of 64 neighbourhood bakeries in a city had an average weekly supply cost of 230 thousand COP. The average weekly supply cost for all neighbourhood bakeries in the city is unknown. The known population standard deviation of individual bakeries' weekly costs is 36 thousand COP. Two reports, at 90\% and 99\% confidence, concern this same sample.

In a separate, completed repeated-sampling study, 250 fresh independent samples of 64 bakeries were drawn from a normal population with a fixed known average cost of 228 thousand COP and the same population standard deviation as the original population. Both unrounded confidence intervals were constructed from each sample. The 90\% intervals missed the known mean 28 times, and the 99\% intervals missed it 4 times.

\textbf{Reporting requirement.} Round reported endpoints to the nearest thousand COP.


\task{1}{Identify and organize the information; calculate SE.}
\BlankInformationTable

\clearpage
\problemtitle{6}{Continued -- interval estimation and coverage}

\task{2}{Calculate MoE and state the interval.}
\EstimationTable{
$\bar{x}$ & Shared point estimate & & & \\[30pt]\midrule
\multicolumn{5}{@{}l}{\textbf{90\% confidence}}\\[3pt]
$\mathrm{MoE}$ & Margin of error & & & \\[30pt]
$L$ & Lower endpoint (unrounded) & & & \\[30pt]
$U$ & Upper endpoint (unrounded) & & & \\[30pt]
$[L,U]$ & Confidence interval (unrounded) & & & \\[30pt]
$W$ & Full width (unrounded) & & & \\[30pt]
$L_{\rm rep}$ & Reported lower endpoint & & & \\[30pt]
$U_{\rm rep}$ & Reported upper endpoint & & & \\[30pt]
$[L_{\rm rep},U_{\rm rep}]$ & Reported confidence interval & & & \\[30pt]
$W_{\rm rep}$ & Reported full width & & & \\[30pt]
}
\clearpage
\problemtitle{6}{Continued -- Task 2}
\EstimationTable{
\multicolumn{5}{@{}l}{\textbf{99\% confidence}}\\[3pt]
$\mathrm{MoE}$ & Margin of error & & & \\[30pt]
$L$ & Lower endpoint (unrounded) & & & \\[30pt]
$U$ & Upper endpoint (unrounded) & & & \\[30pt]
$[L,U]$ & Confidence interval (unrounded) & & & \\[30pt]
$W$ & Full width (unrounded) & & & \\[30pt]
$L_{\rm rep}$ & Reported lower endpoint & & & \\[30pt]
$U_{\rm rep}$ & Reported upper endpoint & & & \\[30pt]
$[L_{\rm rep},U_{\rm rep}]$ & Reported confidence interval & & & \\[30pt]
$W_{\rm rep}$ & Reported full width & & & \\[30pt]
}
\textbf{Compare the centre, widths, and nesting of the intervals.}
Explain how rounding affects the reported widths, while nominal confidence refers to the unrounded procedures.
\workspace{.65in}

\task{3}{Calculate observed coverage and answer the multiple-choice question.}

\choices{The 99\% unrounded procedure has greater long-run coverage and gives a wider interval from this sample; each calculated interval either covers or misses.}{The 99\% interval contains 99\% of individual bakeries' costs.}{The simulation must give exactly 90\% and 99\% coverage.}{The reported 99\% interval has a 99\% probability of containing its fixed mean.}
\vfill

\clearpage
\problemtitle{7}{Clinic waiting times in a simulation}
A random sample of 49 appointments at a clinic in a specified month had an average waiting time of 26 minutes. The average waiting time for all appointments at the clinic that month is unknown. The known population standard deviation of individual appointment waiting times is 14 minutes. The confidence level for the report is 90\%.

In a separate, completed repeated-sampling study, 20 fresh independent samples of 49 appointments were drawn from a normal population with a fixed known average waiting time of 27 minutes and the same population standard deviation as the original population. An unrounded 90\% confidence interval was constructed from each sample. The table groups the counts of intervals containing the known mean by generation order.

\begin{center}\begin{tabular}{@{}lllll@{}}\toprule
Samples in order & 1--5 & 6--10 & 11--15 & 16--20\\\midrule
Covering intervals & 4 & 5 & 4 & 3\\
\bottomrule\end{tabular}\end{center}


\task{1}{Identify and organize the information; calculate SE.}
\BlankInformationTable
\clearpage
\problemtitle{7}{Continued -- interval estimation and coverage}

\task{2}{Calculate MoE and state the interval.}
\EstimationTable{
$\bar{x}$ & Point estimate & & & \\[30pt]\midrule
\multicolumn{5}{@{}l}{\textbf{90\% confidence}}\\[3pt]
$\mathrm{MoE}$ & Margin of error & & & \\[30pt]
$L$ & Lower endpoint & & & \\[30pt]
$U$ & Upper endpoint & & & \\[30pt]
$[L,U]$ & Confidence interval & & & \\[30pt]
$W$ & Full width & & & \\[30pt]
}
\task{3}{Calculate observed coverage and answer the multiple-choice question.}
Is the claim ``Every run of 20 must have exactly 18 covering intervals'' justified?
\choices{Eighteen of 20 cover; every 90\% run must give exactly 18.}{Sixteen of 20 cover; finite runs may differ, while long-run coverage is about 90\% under the model.}{Sixteen of 20 cover; the observed 80\% is the long-run coverage.}{Sixteen of 20 cover; each fixed interval has a 90\% chance of containing the fixed mean.}
\vfill

\clearpage
\problemtitle{8}{Parcel sorting times in a simulation}
A random sample of 64 parcels at a depot during a specified shift had an average sorting time of 18 minutes. The average sorting time for all parcels at the depot during that shift is unknown. The known population standard deviation of individual parcel sorting times is 12 minutes. The confidence level for the report is 90\%.

In a separate, completed repeated-sampling study, 20 fresh independent samples of 64 parcels were drawn from a normal population with a fixed known average sorting time of 19 minutes and the same population standard deviation as the original population. An unrounded 90\% confidence interval was constructed from each sample. The table groups the counts of intervals containing the known mean by generation order.

\begin{center}\begin{tabular}{@{}lllll@{}}\toprule
Samples in order & 1--5 & 6--10 & 11--15 & 16--20\\\midrule
Covering intervals & 5 & 4 & 5 & 5\\
\bottomrule\end{tabular}\end{center}


\task{1}{Identify and organize the information; calculate SE.}
\BlankInformationTable
\clearpage
\problemtitle{8}{Continued -- interval estimation and coverage}

\task{2}{Calculate MoE and state the interval.}
\EstimationTable{
$\bar{x}$ & Point estimate & & & \\[30pt]\midrule
\multicolumn{5}{@{}l}{\textbf{90\% confidence}}\\[3pt]
$\mathrm{MoE}$ & Margin of error & & & \\[30pt]
$L$ & Lower endpoint & & & \\[30pt]
$U$ & Upper endpoint & & & \\[30pt]
$[L,U]$ & Confidence interval & & & \\[30pt]
$W$ & Full width & & & \\[30pt]
}
\task{3}{Calculate observed coverage and answer the multiple-choice question.}
Is the claim ``Every run of 20 must have exactly 18 covering intervals'' justified?
\choices{Nineteen of 20 cover; every 90\% run must give exactly 18, so the procedure is invalid.}{Eighteen of 20 cover; each run must match 90\%.}{Nineteen of 20 cover; 95\% becomes the procedure's long-run confidence.}{Nineteen of 20 cover; the count may vary across finite runs even though long-run coverage is about 90\%.}
\vfill

\clearpage
\problemtitle{9}{Three confidence levels for library processing times}
A random sample of 100 books returned to a library during one month had an average processing time of 24 minutes. The average processing time for all books returned to the library that month is unknown. The known population standard deviation of individual books' processing times is 15 minutes. Three reports, at 90\%, 95\%, and 99\% confidence, concern this same sample.

In a separate, completed repeated-sampling study, 250 fresh independent samples of 100 books were drawn from a normal population with a fixed known average processing time of 25 minutes and the same population standard deviation as the original population. All three unrounded confidence intervals were constructed around each sample's own average. The categories below are mutually exclusive, and each sample appears once.

\textbf{Reporting requirements.} Keep all calculated endpoint digits. Report coverage percentages to two decimal places, in the order 90\%, 95\%, 99\%.

\begin{center}\begin{tabular}{@{}ll@{}}\toprule
Intervals containing the fixed simulation mean & Samples\\\midrule
All three & 219\\
Only 95\% and 99\% & 18\\
Only 99\% & 10\\
None & 3\\
\bottomrule\end{tabular}\end{center}


\task{1}{Identify and organize the information; calculate SE.}
\BlankInformationTable
\clearpage
\problemtitle{9}{Continued -- interval estimation and coverage}

\task{2}{Calculate MoE and state the interval.}
\EstimationTable{
$\bar{x}$ & Shared point estimate & & & \\[30pt]\midrule
\multicolumn{5}{@{}l}{\textbf{90\% confidence}}\\[3pt]
$\mathrm{MoE}$ & Margin of error & & & \\[30pt]
$L$ & Lower endpoint & & & \\[30pt]
$U$ & Upper endpoint & & & \\[30pt]
$[L,U]$ & Confidence interval & & & \\[30pt]
$W$ & Full width & & & \\[30pt]
}
\EstimationTable{
\multicolumn{5}{@{}l}{\textbf{95\% confidence}}\\[3pt]
$\mathrm{MoE}$ & Margin of error & & & \\[30pt]
$L$ & Lower endpoint & & & \\[30pt]
$U$ & Upper endpoint & & & \\[30pt]
$[L,U]$ & Confidence interval & & & \\[30pt]
$W$ & Full width & & & \\[30pt]
}
\clearpage
\problemtitle{9}{Continued -- Task 2}
\EstimationTable{
\multicolumn{5}{@{}l}{\textbf{99\% confidence}}\\[3pt]
$\mathrm{MoE}$ & Margin of error & & & \\[30pt]
$L$ & Lower endpoint & & & \\[30pt]
$U$ & Upper endpoint & & & \\[30pt]
$[L,U]$ & Confidence interval & & & \\[30pt]
$W$ & Full width & & & \\[30pt]
}
\textbf{Compare the centre, widths, and nesting of the intervals.}
\workspace{.65in}

\task{3}{Calculate observed coverage and answer the multiple-choice question.}

\choices{90\%, 95\%, 99\%; every finite run must match its nominal levels.}{87.6\%, 94.8\%, 98.8\%; higher confidence moves the centre closer to the unknown mean.}{87.6\%, 94.8\%, 98.8\%; higher confidence widens intervals about the same sample mean, and describes long-run coverage of the fixed mean.}{87.6\%, 94.8\%, 98.8\%; the calculated 99\% interval has a 99\% probability of containing its fixed mean.}
\vfill

\clearpage
\problemtitle{10}{Three confidence levels for equipment setup times}
A random sample of 81 production runs at a workshop during one month had an average setup time of 62 minutes. The average setup time for all production runs at the workshop that month is unknown. The known population standard deviation of individual runs' setup times is 27 minutes. Three reports, at 90\%, 95\%, and 99\% confidence, concern this same sample.

In a separate, completed repeated-sampling study, 300 fresh independent samples of 81 production runs were drawn from a normal population with a fixed known average setup time of 60 minutes and the same population standard deviation as the original population. All three unrounded confidence intervals were constructed around each sample's own average. The categories below are mutually exclusive, and each sample appears once.

\textbf{Reporting requirements.} Keep all calculated endpoint digits. Report coverage percentages to two decimal places, in the order 90\%, 95\%, 99\%.

\begin{center}\begin{tabular}{@{}ll@{}}\toprule
Intervals containing the fixed simulation mean & Samples\\\midrule
All three & 268\\
Only 95\% and 99\% & 16\\
Only 99\% & 12\\
None & 4\\
\bottomrule\end{tabular}\end{center}


\task{1}{Identify and organize the information; calculate SE.}
\BlankInformationTable
\clearpage
\problemtitle{10}{Continued -- interval estimation and coverage}

\task{2}{Calculate MoE and state the interval.}
\EstimationTable{
$\bar{x}$ & Shared point estimate & & & \\[30pt]\midrule
\multicolumn{5}{@{}l}{\textbf{90\% confidence}}\\[3pt]
$\mathrm{MoE}$ & Margin of error & & & \\[30pt]
$L$ & Lower endpoint & & & \\[30pt]
$U$ & Upper endpoint & & & \\[30pt]
$[L,U]$ & Confidence interval & & & \\[30pt]
$W$ & Full width & & & \\[30pt]
}
\EstimationTable{
\multicolumn{5}{@{}l}{\textbf{95\% confidence}}\\[3pt]
$\mathrm{MoE}$ & Margin of error & & & \\[30pt]
$L$ & Lower endpoint & & & \\[30pt]
$U$ & Upper endpoint & & & \\[30pt]
$[L,U]$ & Confidence interval & & & \\[30pt]
$W$ & Full width & & & \\[30pt]
}
\clearpage
\problemtitle{10}{Continued -- Task 2}
\EstimationTable{
\multicolumn{5}{@{}l}{\textbf{99\% confidence}}\\[3pt]
$\mathrm{MoE}$ & Margin of error & & & \\[30pt]
$L$ & Lower endpoint & & & \\[30pt]
$U$ & Upper endpoint & & & \\[30pt]
$[L,U]$ & Confidence interval & & & \\[30pt]
$W$ & Full width & & & \\[30pt]
}
\textbf{Compare the centre, widths, and nesting of the intervals.}
\workspace{.65in}

\task{3}{Calculate observed coverage and answer the multiple-choice question.}

\choices{About 89.33\%, 94.67\%, 98.67\%; confidence describes long-run coverage, and greater confidence widens intervals with the same centre.}{90\%, 95\%, 99\%; finite runs must match all three levels.}{About 89.33\%, 94.67\%, 98.67\%; higher confidence changes the sample mean to approach the fixed mean.}{About 89.33\%, 94.67\%, 98.67\%; 99\% of individual setup times lie in the calculated 99\% interval.}
\vfill
\end{document}
